% Options for packages loaded elsewhere
\PassOptionsToPackage{unicode}{hyperref}
\PassOptionsToPackage{hyphens}{url}
\documentclass[
  11pt,
]{article}
\usepackage{xcolor}
\usepackage[margin=1in]{geometry}
\usepackage{amsmath,amssymb}
\setcounter{secnumdepth}{-\maxdimen} % remove section numbering
\usepackage{iftex}
\ifPDFTeX
  \usepackage[T1]{fontenc}
  \usepackage[utf8]{inputenc}
  \usepackage{textcomp} % provide euro and other symbols
\else % if luatex or xetex
  \usepackage{unicode-math} % this also loads fontspec
  \defaultfontfeatures{Scale=MatchLowercase}
  \defaultfontfeatures[\rmfamily]{Ligatures=TeX,Scale=1}
\fi
\usepackage{lmodern}
\ifPDFTeX\else
  % xetex/luatex font selection
\fi
% Use upquote if available, for straight quotes in verbatim environments
\IfFileExists{upquote.sty}{\usepackage{upquote}}{}
\IfFileExists{microtype.sty}{% use microtype if available
  \usepackage[]{microtype}
  \UseMicrotypeSet[protrusion]{basicmath} % disable protrusion for tt fonts
}{}
\makeatletter
\@ifundefined{KOMAClassName}{% if non-KOMA class
  \IfFileExists{parskip.sty}{%
    \usepackage{parskip}
  }{% else
    \setlength{\parindent}{0pt}
    \setlength{\parskip}{6pt plus 2pt minus 1pt}}
}{% if KOMA class
  \KOMAoptions{parskip=half}}
\makeatother
\usepackage{color}
\usepackage{fancyvrb}
\newcommand{\VerbBar}{|}
\newcommand{\VERB}{\Verb[commandchars=\\\{\}]}
\DefineVerbatimEnvironment{Highlighting}{Verbatim}{commandchars=\\\{\}}
% Add ',fontsize=\small' for more characters per line
\newenvironment{Shaded}{}{}
\newcommand{\AlertTok}[1]{\textcolor[rgb]{1.00,0.00,0.00}{\textbf{#1}}}
\newcommand{\AnnotationTok}[1]{\textcolor[rgb]{0.38,0.63,0.69}{\textbf{\textit{#1}}}}
\newcommand{\AttributeTok}[1]{\textcolor[rgb]{0.49,0.56,0.16}{#1}}
\newcommand{\BaseNTok}[1]{\textcolor[rgb]{0.25,0.63,0.44}{#1}}
\newcommand{\BuiltInTok}[1]{\textcolor[rgb]{0.00,0.50,0.00}{#1}}
\newcommand{\CharTok}[1]{\textcolor[rgb]{0.25,0.44,0.63}{#1}}
\newcommand{\CommentTok}[1]{\textcolor[rgb]{0.38,0.63,0.69}{\textit{#1}}}
\newcommand{\CommentVarTok}[1]{\textcolor[rgb]{0.38,0.63,0.69}{\textbf{\textit{#1}}}}
\newcommand{\ConstantTok}[1]{\textcolor[rgb]{0.53,0.00,0.00}{#1}}
\newcommand{\ControlFlowTok}[1]{\textcolor[rgb]{0.00,0.44,0.13}{\textbf{#1}}}
\newcommand{\DataTypeTok}[1]{\textcolor[rgb]{0.56,0.13,0.00}{#1}}
\newcommand{\DecValTok}[1]{\textcolor[rgb]{0.25,0.63,0.44}{#1}}
\newcommand{\DocumentationTok}[1]{\textcolor[rgb]{0.73,0.13,0.13}{\textit{#1}}}
\newcommand{\ErrorTok}[1]{\textcolor[rgb]{1.00,0.00,0.00}{\textbf{#1}}}
\newcommand{\ExtensionTok}[1]{#1}
\newcommand{\FloatTok}[1]{\textcolor[rgb]{0.25,0.63,0.44}{#1}}
\newcommand{\FunctionTok}[1]{\textcolor[rgb]{0.02,0.16,0.49}{#1}}
\newcommand{\ImportTok}[1]{\textcolor[rgb]{0.00,0.50,0.00}{\textbf{#1}}}
\newcommand{\InformationTok}[1]{\textcolor[rgb]{0.38,0.63,0.69}{\textbf{\textit{#1}}}}
\newcommand{\KeywordTok}[1]{\textcolor[rgb]{0.00,0.44,0.13}{\textbf{#1}}}
\newcommand{\NormalTok}[1]{#1}
\newcommand{\OperatorTok}[1]{\textcolor[rgb]{0.40,0.40,0.40}{#1}}
\newcommand{\OtherTok}[1]{\textcolor[rgb]{0.00,0.44,0.13}{#1}}
\newcommand{\PreprocessorTok}[1]{\textcolor[rgb]{0.74,0.48,0.00}{#1}}
\newcommand{\RegionMarkerTok}[1]{#1}
\newcommand{\SpecialCharTok}[1]{\textcolor[rgb]{0.25,0.44,0.63}{#1}}
\newcommand{\SpecialStringTok}[1]{\textcolor[rgb]{0.73,0.40,0.53}{#1}}
\newcommand{\StringTok}[1]{\textcolor[rgb]{0.25,0.44,0.63}{#1}}
\newcommand{\VariableTok}[1]{\textcolor[rgb]{0.10,0.09,0.49}{#1}}
\newcommand{\VerbatimStringTok}[1]{\textcolor[rgb]{0.25,0.44,0.63}{#1}}
\newcommand{\WarningTok}[1]{\textcolor[rgb]{0.38,0.63,0.69}{\textbf{\textit{#1}}}}
\setlength{\emergencystretch}{3em} % prevent overfull lines
\providecommand{\tightlist}{%
  \setlength{\itemsep}{0pt}\setlength{\parskip}{0pt}}
\usepackage{bookmark}
\IfFileExists{xurl.sty}{\usepackage{xurl}}{} % add URL line breaks if available
\urlstyle{same}
\hypersetup{
  hidelinks,
  pdfcreator={LaTeX via pandoc}}

\author{}
\date{}

\begin{document}

\section{Reference Material}\label{reference-material}

\subsection{Iterator Methods and
Adapters}\label{iterator-methods-and-adapters}

\begin{itemize}
\tightlist
\item
  \texttt{next(\&mut\ self)\ -\textgreater{}\ Option\textless{}Self::Item\textgreater{}}
  Get the next element of an iterator (\texttt{None} if there isn't one)
\item
  \texttt{collect\textless{}B\textgreater{}(self)\ -\textgreater{}\ B}
  -\textgreater{} Transforms the iterator into a collection of type
  \texttt{B}
\item
  \texttt{take(n:\ usize)\ -\textgreater{}\ Take\textless{}Self\textgreater{}}
  -\textgreater{} Take first N elements of an iterator and \emph{turn
  them into an iterator}
\item
  \texttt{skip(n:\ usize)\ -\textgreater{}\ Skip\textless{}Self\textgreater{}}
  -\textgreater{} Skip first N elements of an iterator and \emph{turn
  them into an iterator}
\item
  \texttt{cycle()\ -\textgreater{}\ Cycle\textless{}Self\textgreater{}}
  -\textgreater{} Turn a finite iterator into an infinite one that
  repeats itself
\item
  \texttt{for\_each(\textbar{}x\textbar{}\ ...)\ -\textgreater{}\ ()}
  -\textgreater{} Apply a closure to each element in the iterator
\item
  \texttt{filter(\textbar{}x\textbar{}\ ...)\ -\textgreater{}\ Filter\textless{}Self,\ P\textgreater{}}
  -\textgreater{} \emph{Create new iterator} from old one for elements
  where closure is true
\item
  \texttt{map(\textbar{}x\textbar{}\ ...)\ -\textgreater{}\ Map\textless{}Self,\ F\textgreater{}}
  -\textgreater{} \emph{Create new iterator} by applying closure to
  input iterator
\item
  \texttt{any(\textbar{}x\textbar{}\ ...)\ -\textgreater{}\ bool}
  -\textgreater{} Return \texttt{true} if closure is true for any
  element of the iterator
\item
  \texttt{fold(init:\ B,\ \textbar{}acc,\ x\textbar{}\ ...)\ -\textgreater{}\ B}
  -\textgreater{} Initialize accumulator to \texttt{init}, execute
  closure on each element
\item
  \texttt{reduce(\textbar{}acc,\ x\textbar{}\ ...)\ -\textgreater{}\ Option\textless{}Self::Item\textgreater{}}
  -\textgreater{} Similar to fold but the initial value is the first
  element
\item
  \texttt{zip(other\_iter)\ -\textgreater{}\ Zip\textless{}Self,\ U\textgreater{}}
  -\textgreater{} Zip two iterators together into pairs of
  \texttt{(a,\ b)}
\item
  \texttt{enumerate()\ -\textgreater{}\ Enumerate\textless{}Self\textgreater{}}
  -\textgreater{} Create iterator of \texttt{(index,\ element)} pairs
\item
  \texttt{partition(\textbar{}x\textbar{}\ ...)\ -\textgreater{}\ (B,\ B)}
  -\textgreater{} Split iterator into two \emph{collections} based on
  predicate
\item
  \texttt{sum()\ -\textgreater{}\ S} -\textgreater{} Sum all elements
  (works on numeric iterators)
\item
  \texttt{count()\ -\textgreater{}\ usize} -\textgreater{} Count the
  number of elements in the iterator
\item
  \texttt{copied()\ -\textgreater{}\ Copied\textless{}Self\textgreater{}}
  -\textgreater{} Create iterator of owned copies from an iterator of
  references
\end{itemize}

\subsection{String and \&str Methods}\label{string-and-str-methods}

\begin{itemize}
\tightlist
\item
  \texttt{s.len()\ -\textgreater{}\ usize} -\textgreater{} Number of
  bytes in the string
\item
  \texttt{s.contains("sub")\ -\textgreater{}\ bool} -\textgreater{}
  \texttt{true} if string contains the substring
\item
  \texttt{s.to\_uppercase()\ -\textgreater{}\ String} -\textgreater{}
  Returns a new \texttt{String} in uppercase
\item
  \texttt{s.to\_lowercase()\ -\textgreater{}\ String} -\textgreater{}
  Returns a new \texttt{String} in lowercase
\item
  \texttt{s.push(\textquotesingle{}c\textquotesingle{})\ -\textgreater{}\ ()}
  -\textgreater{} Append a single character (mutates \texttt{String})
\item
  \texttt{s.push\_str("text")\ -\textgreater{}\ ()} -\textgreater{}
  Append a string slice (mutates \texttt{String})
\item
  \texttt{s.chars()\ -\textgreater{}\ Chars\textless{}\textquotesingle{}\_\textgreater{}}
  -\textgreater{} Iterator over the characters of the string
\item
  \texttt{s.split\_whitespace()\ -\textgreater{}\ SplitWhitespace\textless{}\textquotesingle{}\_\textgreater{}}
  -\textgreater{} Iterator of substrings split by whitespace
\item
  \texttt{s.parse::\textless{}T\textgreater{}()\ -\textgreater{}\ Result\textless{}T,\ T::Err\textgreater{}}
  -\textgreater{} Parse the string into type \texttt{T}
\item
  \texttt{s.is\_empty()\ -\textgreater{}\ bool} -\textgreater{}
  \texttt{true} if the string has length 0
\end{itemize}

\subsection{Vec Methods}\label{vec-methods}

\begin{itemize}
\tightlist
\item
  \texttt{v.push(val)\ -\textgreater{}\ ()} -\textgreater{} Append
  element to the end
\item
  \texttt{v.pop()\ -\textgreater{}\ Option\textless{}T\textgreater{}}
  -\textgreater{} Remove and return last element
\item
  \texttt{v.get(i:\ usize)\ -\textgreater{}\ Option\textless{}\&T\textgreater{}}
  -\textgreater{} Safe bounds-checked access
\item
  \texttt{v.len()\ -\textgreater{}\ usize} -\textgreater{} Number of
  elements
\item
  \texttt{v.is\_empty()\ -\textgreater{}\ bool} -\textgreater{}
  \texttt{true} if the vector has no elements
\item
  \texttt{v.iter()\ -\textgreater{}\ Iter\textless{}\textquotesingle{}\_,\ T\textgreater{}}
  -\textgreater{} Iterator of immutable references \texttt{\&T}
\item
  \texttt{v.iter\_mut()\ -\textgreater{}\ IterMut\textless{}\textquotesingle{}\_,\ T\textgreater{}}
  -\textgreater{} Iterator of mutable references \texttt{\&mut\ T}
\item
  \texttt{v.into\_iter()\ -\textgreater{}\ IntoIter\textless{}T\textgreater{}}
  -\textgreater{} Consuming iterator of owned values \texttt{T}
\item
  \texttt{v.sort\_by(\textbar{}a,\ b\textbar{}\ ...)\ -\textgreater{}\ ()}
  -\textgreater{} Sort in place using a comparison closure (closure
  returns \texttt{Ordering})
\end{itemize}

\subsection{Option Methods}\label{option-methods}

\begin{itemize}
\tightlist
\item
  \texttt{opt.unwrap()\ -\textgreater{}\ T} -\textgreater{} Extract the
  value or \textbf{panic} if \texttt{None}
\item
  \texttt{opt.unwrap\_or(default:\ T)\ -\textgreater{}\ T}
  -\textgreater{} Extract the value or return \texttt{default} (eager)
\item
  \texttt{opt.unwrap\_or\_else(\textbar{}\textbar{}\ ...)\ -\textgreater{}\ T}
  -\textgreater{} Extract the value or compute default via closure
  (lazy)
\item
  \texttt{opt.map(\textbar{}x\textbar{}\ ...)\ -\textgreater{}\ Option\textless{}U\textgreater{}}
  -\textgreater{} Transform the inner value; \texttt{None} stays
  \texttt{None}
\item
  \texttt{opt.is\_some()\ -\textgreater{}\ bool} -\textgreater{}
  \texttt{true} if \texttt{Some}
\item
  \texttt{opt.is\_none()\ -\textgreater{}\ bool} -\textgreater{}
  \texttt{true} if \texttt{None}
\end{itemize}

\subsection{Other Useful Methods}\label{other-useful-methods}

\begin{itemize}
\tightlist
\item
  \texttt{val.clone()\ -\textgreater{}\ T} -\textgreater{} Create a deep
  copy of the value
\item
  \texttt{a.cmp(\&b)\ -\textgreater{}\ std::cmp::Ordering}
  -\textgreater{} Compare two values
\item
  \texttt{(x\ as\ f64).sqrt()\ -\textgreater{}\ f64} -\textgreater{}
  Square root of a floating-point number
\item
  \texttt{format!("...",\ args)\ -\textgreater{}\ String}
  -\textgreater{} Like \texttt{println!} but returns a \texttt{String}
  instead of printing
\item
  \texttt{\{:?\}} -\textgreater{} Debug format specifier (in
  \texttt{println!} / \texttt{format!})
\item
  \texttt{\{:.2\}} -\textgreater{} Display with 2 decimal places
\end{itemize}

\subsection{Test Assertion Macros}\label{test-assertion-macros}

\begin{itemize}
\tightlist
\item
  \texttt{assert!(expr)\ -\textgreater{}\ ()} -\textgreater{} Pass if
  \texttt{expr} is \texttt{true}
\item
  \texttt{assert\_eq!(a,\ b)\ -\textgreater{}\ ()} -\textgreater{} Pass
  if \texttt{a\ ==\ b}; shows both values on failure
\item
  \texttt{assert\_ne!(a,\ b)\ -\textgreater{}\ ()} -\textgreater{} Pass
  if \texttt{a\ !=\ b}; shows both values on failure
\item
  Assert macros can take an optional custom failure message, e.g.,
  \texttt{assert\_eq!(a,\ b,\ "\{\}\ is\ not\ equal\ to\ \{\}",\ a,\ b)}.
\end{itemize}

\subsection{HashMap Methods}\label{hashmap-methods}

\subsubsection{Basic Operations:}\label{basic-operations}

\begin{itemize}
\tightlist
\item
  \texttt{new()}: Creates an empty HashMap.
\item
  \texttt{insert(key,\ value)\ -\textgreater{}\ Option\textless{}V\textgreater{}}:
  Adds a key-value pair to the map. If the map did not have the key then
  \texttt{None} is returned, if the map did have the key \emph{the old
  value is returned} as \texttt{Some(V)}.
\item
  \texttt{remove(key)\ -\textgreater{}\ Option\textless{}V\textgreater{}}:
  Removes a key-value pair from the map. If the map did not have the key
  then \texttt{None} is returned, if the map did have the key the value
  is returned as \texttt{Some(V)}.
\item
  \texttt{get(key)\ -\textgreater{}\ Option\textless{}\&V\textgreater{}}:
  Returns a reference to the value in the map, if any, that is equal to
  the given key. Returns \texttt{None} if the key is not present.
\item
  \texttt{contains\_key(key)\ -\textgreater{}\ bool}: Checks if the map
  contains a specific key. Returns \texttt{true} if present,
  \texttt{false} otherwise.
\item
  \texttt{len()\ -\textgreater{}\ usize}: Returns the number of
  key-value pairs in the map.
\item
  \texttt{is\_empty()\ -\textgreater{}\ bool}: Checks if the map
  contains no key-value pairs.
\item
  \texttt{clear()}: Removes all key-value pairs from the map.
\item
  \texttt{drain()}: Clears the map and returns all key-value pairs as an
  iterator.
\end{itemize}

\subsubsection{Entry API:}\label{entry-api}

\begin{itemize}
\tightlist
\item
  \texttt{entry(key).or\_insert(value)}: Returns a mutable reference to
  the value in the map, if any, that is equal to the given key. If the
  key is not present, inserts the given value and returns a mutable
  reference to the new value.
\end{itemize}

\subsubsection{Iterators and Views:}\label{iterators-and-views}

\begin{itemize}
\tightlist
\item
  \texttt{iter()}: Returns an immutable iterator over the key-value
  pairs in the map.
\item
  \texttt{iter\_mut()}: Returns a mutable iterator over the key-value
  pairs in the map.
\item
  \texttt{keys()}: Returns an iterator over the keys in the map.
\item
  \texttt{values()}: Returns an iterator over the values in the map.
\item
  \texttt{values\_mut()}: Returns a mutable iterator over the values in
  the map.
\end{itemize}

\subsection{HashSet Methods}\label{hashset-methods}

\subsubsection{Basic Operations:}\label{basic-operations-1}

\begin{itemize}
\tightlist
\item
  \texttt{new()}: Creates an empty HashSet.
\item
  \texttt{insert(value)}: Adds a value to the set. Returns true if the
  value was not present, false otherwise.
\item
  \texttt{remove(value)}: Removes a value from the set. Returns true if
  the value was present, false otherwise.
\item
  \texttt{contains(value)}: Checks if the set contains a specific value.
  Returns true if present, false otherwise.
\item
  \texttt{len()}: Returns the number of elements in the set.
\item
  \texttt{is\_empty()}: Checks if the set contains no elements.
\item
  \texttt{clear()}: Removes all elements from the set.
\item
  \texttt{drain()}: Returns an iterator that removes all elements and
  yields them. The set becomes empty after this operation.
\end{itemize}

\subsubsection{Entry API:}\label{entry-api-1}

\begin{itemize}
\tightlist
\item
  \texttt{entry(value).or\_insert(value)}: Returns a mutable reference
  to the value in the set, if any, that is equal to the given value. If
  the value is not present, inserts the given value and returns a
  mutable reference to the new value.
\end{itemize}

\subsubsection{Set Operations:}\label{set-operations}

\begin{itemize}
\tightlist
\item
  \texttt{union(\&self,\ other:\ \&HashSet\textless{}T\textgreater{})}:
  Returns an iterator over the elements that are in self or other (or
  both).
\item
  \texttt{intersection(\&self,\ other:\ \&HashSet\textless{}T\textgreater{})}:
  Returns an iterator over the elements that are in both self and other.
\item
  \texttt{difference(\&self,\ other:\ \&HashSet\textless{}T\textgreater{})}:
  Returns an iterator over the elements that are in self but not in
  other.
\item
  \texttt{symmetric\_difference(\&self,\ other:\ \&HashSet\textless{}T\textgreater{})}:
  Returns an iterator over the elements that are in self or other, but
  not in both.
\item
  \texttt{is\_subset(\&self,\ other:\ \&HashSet\textless{}T\textgreater{})}:
  Checks if self is a subset of other.
\item
  \texttt{is\_superset(\&self,\ other:\ \&HashSet\textless{}T\textgreater{})}:
  Checks if self is a superset of other.
\item
  \texttt{is\_disjoint(\&self,\ other:\ \&HashSet\textless{}T\textgreater{})}:
  Checks if self has no elements in common with other.
\end{itemize}

\subsubsection{Iterators and Views:}\label{iterators-and-views-1}

\begin{itemize}
\tightlist
\item
  \texttt{iter()}: Returns an immutable iterator over the elements in
  the set.
\item
  \texttt{get(value)}: Returns a reference to the value in the set, if
  any, that is equal to the given value.
\end{itemize}

\subsection{BTreeMap Methods}\label{btreemap-methods}

\subsubsection{Basic Operations:}\label{basic-operations-2}

\begin{itemize}
\tightlist
\item
  \texttt{new()}: Creates an empty BTreeMap.
\item
  \texttt{insert(key,\ value)\ -\textgreater{}\ Option\textless{}V\textgreater{}}:
  Adds a key-value pair to the map. If the map did not have the key then
  \texttt{None} is returned, if the map did have the key \emph{the old
  value is returned} as \texttt{Some(V)}.
\item
  \texttt{remove(key)\ -\textgreater{}\ Option\textless{}V\textgreater{}}:
  Removes a key-value pair from the map. If the map did not have the key
  then \texttt{None} is returned, if the map did have the key the value
  is returned as \texttt{Some(V)}.
\item
  \texttt{get(key)\ -\textgreater{}\ Option\textless{}\&V\textgreater{}}:
  Returns a reference to the value in the map, if any, that is equal to
  the given key. Returns \texttt{None} if the key is not present.
\item
  \texttt{contains\_key(key)\ -\textgreater{}\ bool}: Checks if the map
  contains a specific key. Returns \texttt{true} if present,
  \texttt{false} otherwise.
\item
  \texttt{len()\ -\textgreater{}\ usize}: Returns the number of
  key-value pairs in the map.
\item
  \texttt{is\_empty()\ -\textgreater{}\ bool}: Checks if the map
  contains no key-value pairs.
\item
  \texttt{clear()}: Removes all key-value pairs from the map.
\item
  \texttt{drain()}: Clears the map and returns all key-value pairs as an
  iterator.
\end{itemize}

\subsubsection{Entry API:}\label{entry-api-2}

\begin{itemize}
\tightlist
\item
  \texttt{entry(key).or\_insert(value)}: Returns a mutable reference to
  the value in the map, if any, that is equal to the given key. If the
  key is not present, inserts the given value and returns a mutable
  reference to the new value.
\end{itemize}

\subsubsection{Iterators and Views:}\label{iterators-and-views-2}

\begin{itemize}
\tightlist
\item
  \texttt{iter()}: Returns an immutable iterator over the key-value
  pairs in the map.
\item
  \texttt{iter\_mut()}: Returns a mutable iterator over the key-value
  pairs in the map.
\item
  \texttt{keys()}: Returns an iterator over the keys in the map.
\item
  \texttt{values()}: Returns an iterator over the values in the map.
\item
  \texttt{values\_mut()}: Returns a mutable iterator over the values in
  the map.
\item
  \texttt{range(start..end)}: Returns an iterator over the key-value
  pairs in the map in the range \texttt{{[}start,\ end)}.
\item
  \texttt{first\_key\_value()}: Returns the first key-value pair in the
  map.
\item
  \texttt{last\_key\_value()}: Returns the last key-value pair in the
  map.
\end{itemize}

\subsection{BTreeSet Methods}\label{btreeset-methods}

\subsubsection{Basic Operations:}\label{basic-operations-3}

\begin{itemize}
\tightlist
\item
  \texttt{new()}: Creates an empty BTreeSet.
\item
  \texttt{insert(value)}: Adds a value to the set. Returns true if the
  value was not present, false otherwise.
\item
  \texttt{remove(value)}: Removes a value from the set. Returns true if
  the value was present, false otherwise.
\item
  \texttt{contains(value)}: Checks if the set contains a specific value.
  Returns true if present, false otherwise.
\item
  \texttt{len()}: Returns the number of elements in the set.
\item
  \texttt{is\_empty()}: Checks if the set contains no elements.
\item
  \texttt{clear()}: Removes all elements from the set.
\item
  \texttt{drain()}: Returns an iterator that removes all elements and
  yields them. The set becomes empty after this operation.
\end{itemize}

\subsubsection{Entry API:}\label{entry-api-3}

\begin{itemize}
\tightlist
\item
  \texttt{entry(value).or\_insert(value)}: Returns a mutable reference
  to the value in the set, if any, that is equal to the given value. If
  the value is not present, inserts the given value and returns a
  mutable reference to the new value.
\end{itemize}

\subsubsection{Set Operations:}\label{set-operations-1}

\begin{itemize}
\tightlist
\item
  \texttt{union(\&self,\ other:\ \&BTreeSet\textless{}T\textgreater{})}:
  Returns an iterator over the elements that are in self or other (or
  both).
\item
  \texttt{intersection(\&self,\ other:\ \&BTreeSet\textless{}T\textgreater{})}:
  Returns an iterator over the elements that are in both self and other.
\item
  \texttt{difference(\&self,\ other:\ \&BTreeSet\textless{}T\textgreater{})}:
  Returns an iterator over the elements that are in self but not in
  other.
\item
  \texttt{symmetric\_difference(\&self,\ other:\ \&BTreeSet\textless{}T\textgreater{})}:
  Returns an iterator over the elements that are in self or other, but
  not in both.
\item
  \texttt{is\_subset(\&self,\ other:\ \&BTreeSet\textless{}T\textgreater{})}:
  Checks if self is a subset of other.
\item
  \texttt{is\_superset(\&self,\ other:\ \&BTreeSet\textless{}T\textgreater{})}:
  Checks if self is a superset of other.
\item
  \texttt{is\_disjoint(\&self,\ other:\ \&BTreeSet\textless{}T\textgreater{})}:
  Checks if self has no elements in common with other.
\end{itemize}

\subsubsection{Iterators and Views:}\label{iterators-and-views-3}

\begin{itemize}
\tightlist
\item
  \texttt{iter()}: Returns an immutable iterator over the elements in
  the set.
\item
  \texttt{get(value)}: Returns a reference to the value in the set, if
  any, that is equal to the given value.
\end{itemize}

\subsection{BinaryHeap Methods
(Max-Heap)}\label{binaryheap-methods-max-heap}

\texttt{std::collections::BinaryHeap} is a max-heap priority queue. Use
\texttt{std::cmp::Reverse} to get a min-heap.

\begin{itemize}
\tightlist
\item
  \texttt{new()}: Creates an empty \texttt{BinaryHeap}.
\item
  \texttt{from(vals)}: Creates a \texttt{BinaryHeap} from an array or
  vector.
\item
  \texttt{from\_iter(iter)}: Creates a \texttt{BinaryHeap} from an
  iterator.
\item
  \texttt{push(value)\ -\textgreater{}\ ()}: Inserts a value into the
  heap.
\item
  \texttt{pop()\ -\textgreater{}\ Option\textless{}T\textgreater{}}:
  Removes and returns the largest element, or \texttt{None} if empty.
\item
  \texttt{peek()\ -\textgreater{}\ Option\textless{}\&T\textgreater{}}:
  Returns a reference to the largest element without removing it.
\item
  \texttt{len()\ -\textgreater{}\ usize}: Number of elements.
\item
  \texttt{is\_empty()\ -\textgreater{}\ bool}: \texttt{true} if the heap
  has no elements.
\item
  \texttt{iter()}: Returns an iterator over the elements in arbitrary
  order (not sorted).
\item
  \texttt{into\_sorted\_vec()\ -\textgreater{}\ Vec\textless{}T\textgreater{}}:
  Consumes the heap and returns a \texttt{Vec} sorted in ascending
  order.
\end{itemize}

\subsection{VecDeque Methods (Double-Ended
Queue)}\label{vecdeque-methods-double-ended-queue}

\texttt{std::collections::VecDeque} is a double-ended queue (ring
buffer). Commonly used as a FIFO queue via \texttt{push\_back} /
\texttt{pop\_front}.

\begin{itemize}
\tightlist
\item
  \texttt{new()}: Creates an empty \texttt{VecDeque}.
\item
  \texttt{push\_back(value)\ -\textgreater{}\ ()}: Appends an element to
  the back.
\item
  \texttt{push\_front(value)\ -\textgreater{}\ ()}: Prepends an element
  to the front.
\item
  \texttt{pop\_back()\ -\textgreater{}\ Option\textless{}T\textgreater{}}:
  Removes and returns the back element.
\item
  \texttt{pop\_front()\ -\textgreater{}\ Option\textless{}T\textgreater{}}:
  Removes and returns the front element.
\item
  \texttt{front()\ -\textgreater{}\ Option\textless{}\&T\textgreater{}}:
  Reference to the front element.
\item
  \texttt{back()\ -\textgreater{}\ Option\textless{}\&T\textgreater{}}:
  Reference to the back element.
\item
  \texttt{len()\ -\textgreater{}\ usize}: Number of elements.
\item
  \texttt{is\_empty()\ -\textgreater{}\ bool}: \texttt{true} if empty.
\item
  \texttt{iter()}: Iterator from front to back.
\end{itemize}

\subsection{LinkedList Methods}\label{linkedlist-methods}

\texttt{std::collections::LinkedList} is a doubly-linked list.

\begin{itemize}
\tightlist
\item
  \texttt{new()}: Creates an empty \texttt{LinkedList}.
\item
  \texttt{push\_back(value)\ -\textgreater{}\ ()}: Appends an element to
  the back.
\item
  \texttt{push\_front(value)\ -\textgreater{}\ ()}: Prepends an element
  to the front.
\item
  \texttt{pop\_back()\ -\textgreater{}\ Option\textless{}T\textgreater{}}:
  Removes and returns the back element.
\item
  \texttt{pop\_front()\ -\textgreater{}\ Option\textless{}T\textgreater{}}:
  Removes and returns the front element.
\item
  \texttt{front()\ -\textgreater{}\ Option\textless{}\&T\textgreater{}}:
  Reference to the front element.
\item
  \texttt{back()\ -\textgreater{}\ Option\textless{}\&T\textgreater{}}:
  Reference to the back element.
\item
  \texttt{len()\ -\textgreater{}\ usize}: Number of elements.
\item
  \texttt{is\_empty()\ -\textgreater{}\ bool}: \texttt{true} if empty.
\item
  \texttt{iter()}: Iterator from front to back.
\end{itemize}

\subsection{Common Derivable Traits}\label{common-derivable-traits}

Add with \texttt{\#{[}derive(...){]}} on a \texttt{struct} or
\texttt{enum}:

\begin{itemize}
\tightlist
\item
  \texttt{Debug}: Enables formatting with \texttt{\{:?\}} in
  \texttt{println!} / \texttt{format!}.
\item
  \texttt{Clone}: Provides \texttt{.clone()} to produce a deep copy.
\item
  \texttt{Copy}: Value is copied (bitwise) on assignment instead of
  moved. Requires \texttt{Clone}. Only valid for types whose fields are
  all \texttt{Copy} (e.g., no \texttt{String}, no
  \texttt{Vec\textless{}T\textgreater{}}).
\item
  \texttt{PartialEq}: Enables \texttt{==} and \texttt{!=}.
\item
  \texttt{Eq}: Marker trait for total equality (no \texttt{NaN}-like
  values). Requires \texttt{PartialEq}. Needed (with \texttt{Hash}) for
  \texttt{HashMap} / \texttt{HashSet} keys.
\item
  \texttt{PartialOrd}: Enables \texttt{\textless{}},
  \texttt{\textless{}=}, \texttt{\textgreater{}},
  \texttt{\textgreater{}=} via \texttt{partial\_cmp}.
\item
  \texttt{Ord}: Total ordering via \texttt{cmp}. Requires \texttt{Eq}
  and \texttt{PartialOrd}. Needed for \texttt{BTreeMap} /
  \texttt{BTreeSet} keys and \texttt{BinaryHeap} elements.
\item
  \texttt{Hash}: Enables use as a \texttt{HashMap} / \texttt{HashSet}
  key. Typically combined with \texttt{Eq}.
\item
  \texttt{Default}: Provides \texttt{T::default()} that returns a
  default value.
\end{itemize}

Example:

\begin{Shaded}
\begin{Highlighting}[]
\AttributeTok{\#[}\NormalTok{derive}\AttributeTok{(}\BuiltInTok{Debug}\OperatorTok{,} \BuiltInTok{Clone}\OperatorTok{,} \BuiltInTok{PartialEq}\OperatorTok{,} \BuiltInTok{Eq}\OperatorTok{,} \BuiltInTok{Hash}\AttributeTok{)]}
\KeywordTok{struct}\NormalTok{ Point }\OperatorTok{\{}\NormalTok{ x}\OperatorTok{:} \DataTypeTok{i32}\OperatorTok{,}\NormalTok{ y}\OperatorTok{:} \DataTypeTok{i32} \OperatorTok{\}}
\end{Highlighting}
\end{Shaded}


\end{document}
